\begin{ouverture}[calcul]{6}{Comparer et calculer des fractions}{Partie 1 · Nombres et calculs}
{Égalités · Comparaisons · Sommes · Produit par un entier}
\objectif{Établir des égalités de fractions et simplifier.}
\objectif{Comparer, encadrer et ranger des fractions et des nombres mixtes.}
\objectif{Additionner, soustraire, multiplier par un entier et appliquer une fraction.}
\prerequis{Fraction quotient ; tables, multiples et diviseurs.}
\end{ouverture}
\pagestyle{otMain}\essentiel\fontsize{11}{15}\selectfont
\phantomsection\label{cours-6-1}
\essSection{1}{Fractions égales et simplification}
\begin{propriete}On ne change pas une fraction en multipliant, ou en divisant, son numérateur et son dénominateur par un même nombre non nul : $\frac23=\frac{2\times4}{3\times4}=\frac8{12}$ ; $\frac{18}{24}=\frac{18:6}{24:6}=\frac34$.\end{propriete}
\label{fig-cours-6-1-cours-bandes-egales}\figureGeo{cours-bandes-egales}

\begin{savoirfaire}[Simplifier]On cherche un diviseur commun grâce aux tables : 18 et 24 sont dans la table de 6. On peut aussi simplifier en plusieurs fois : $\frac{18}{24}=\frac9{12}=\frac34$.\end{savoirfaire}
\essSection{2}{Comparer, encadrer et ranger}
\begin{propriete}Avec le même dénominateur, la plus grande fraction a le plus grand numérateur. Sinon, on se ramène au même dénominateur : $\frac56=\frac{10}{12}>\frac7{12}$. Avec le même numérateur, la plus grande a le plus petit dénominateur : $\frac45>\frac49$.\end{propriete}
\begin{exemple}Pour ranger $\frac74$, $\frac32$ et $1+\frac18$, on écrit tout en huitièmes : $\frac{14}8$, $\frac{12}8$ et $\frac98$. Donc $1+\frac18<\frac32<\frac74$.\end{exemple}
\clearpage
\phantomsection\label{cours-6-2}
\essSection{3}{Additionner et soustraire des fractions}
\begin{propriete}Pour additionner ou soustraire des fractions, on les écrit avec le même dénominateur, puis on ajoute ou on retranche les numérateurs : $\frac34+\frac18=\frac68+\frac18=\frac78$ ; $\frac23+\frac14=\frac8{12}+\frac3{12}=\frac{11}{12}$.\end{propriete}
\begin{exemple}À savoir sans calcul : $\frac12+\frac12=1$ ; $\frac14+\frac14=\frac12$ ; $1-\frac14=\frac34$ ; $\frac12+\frac14=\frac34$ ; $1-\frac12=\frac12$ ; $\frac34+\frac14=1$ ; $\frac12-\frac14=\frac14$ ; $\frac34-\frac14=\frac12$.\end{exemple}
\begin{attention}On n’additionne jamais les dénominateurs : $\frac14+\frac14$ vaut $\frac24$, deux quarts, et non $\frac28$.\end{attention}
\essSection{4}{Multiplier par un entier, appliquer une fraction}
\begin{propriete}$7\times\frac14=\frac74$ : sept fois un quart. C’est aussi $\frac14\times7$, le quart de 7. De même $5\times\frac23=\frac{10}3$ : on multiplie le numérateur par l’entier.\end{propriete}
\begin{savoirfaire}[Prendre une fraction d’un nombre]$\frac23$ de 12, c’est $\frac23\times12$ : $12:3\times2=8$. De même, $\frac34$ de 10 m vaut $10:4\times3=7{,}5$ m.\end{savoirfaire}

\clearpage\rubrique{Découverte}
\begin{blocLivre}{ex-6-1}\begin{exercice}[1]{Fractions égales}
Compléter : $\frac23=\frac{\square}{12}$ ; $\frac54=\frac{20}{\square}$ ; $\frac{18}{24}=\frac3{\square}$.
\end{exercice}\end{blocLivre}
\begin{blocLivre}{ex-6-2}\begin{exercice}[1]{Simplifier}
Simplifier le plus possible : $\frac{15}{20}$ ; $\frac{24}{36}$ ; $\frac{45}{10}$.
\end{exercice}\end{blocLivre}
\begin{blocLivre}{ex-6-3}\begin{exercice}[1]{Comparer}
Comparer avec $<$ ou $>$ : $\frac57$ et $\frac37$ ; $\frac56$ et $\frac7{12}$ ; $\frac45$ et $\frac49$ ; $\frac98$ et 1.
\end{exercice}\end{blocLivre}
\begin{blocLivre}{ex-6-4}\begin{exercice}[1]{Calcul mental}
Calculer : $\frac12+\frac14$ ; $1-\frac14$ ; $\frac34-\frac12$ ; $\frac14+\frac14+\frac12$.
\end{exercice}\end{blocLivre}
\begin{blocLivre}{ex-6-5}\begin{exercice}[1]{Fraction d’un nombre}
Calculer : $\frac23$ de 12 ; $\frac34$ de 10 ; $\frac15$ de 35.
\end{exercice}\end{blocLivre}

\rubrique{Entraînement}
\begin{blocLivre}{ex-6-6}\begin{exercice}[2]{Ranger}
Ranger dans l’ordre croissant : $\frac74$ ; $\frac32$ ; $1+\frac18$ ; $\frac{15}8$.
\end{exercice}\end{blocLivre}
\begin{blocLivre}{ex-6-7}\begin{exercice}[2]{Sommes et différences}
Calculer : $\frac34+\frac58$ ; $\frac{11}{12}-\frac13$ ; $\frac23+\frac14$.
\end{exercice}\end{blocLivre}
\begin{blocLivre}{ex-6-8}\begin{exercice}[2]{Multiplier par un entier}
Calculer : $7\times\frac14$ ; $5\times\frac23$ ; $6\times\frac56$ ; $\frac14\times7$.
\end{exercice}\end{blocLivre}
\begin{blocLivre}{ex-6-9}\begin{exercice}[2]{Sur une même droite}
L’unité est partagée en douzièmes. Écrire les abscisses de A, B et C en douzièmes, les simplifier, puis ranger les trois nombres.
\label{fig-ex-6-9}\figureGeo{fc-9}
\end{exercice}\end{blocLivre}
\begin{blocLivre}{ex-6-10}\begin{exercice}[2]{Égales ou pas ?}
Les fractions $\frac69$ et $\frac{10}{15}$ sont-elles égales ? Et $\frac46$ et $\frac68$ ?
\end{exercice}\end{blocLivre}

\rubrique{Problèmes}
\begin{blocLivre}{ex-6-11}\begin{exercice}[3]{Le gâteau}
Emma mange le quart d’un gâteau et Lucas le tiers. Quelle fraction du gâteau ont-ils mangée ensemble ? Quelle fraction reste-t-il ?
\end{exercice}\end{blocLivre}
\begin{blocLivre}{ex-6-12}\begin{exercice}[3]{La cantine}
Dans une classe de 28 élèves, les trois quarts déjeunent à la cantine. Combien d’élèves déjeunent à la cantine ? Quelle fraction de la classe ne s’y rend pas, et combien d’élèves cela représente-t-il ?
\end{exercice}\end{blocLivre}
\begin{blocLivre}{ex-6-13}\begin{exercice}[3]{La recette}
Une recette pour 4 personnes demande $\frac34$ de litre de lait. Quelle quantité faut-il pour 12 personnes ? Donner le résultat en fraction, puis en écriture décimale.
\end{exercice}\end{blocLivre}
\begin{blocLivre}{ex-6-14}\begin{exercice}[3]{Le réservoir}
Un réservoir de 60 L est rempli aux $\frac56$. On utilise ensuite une quantité égale au quart du réservoir plein. Combien de litres restent ? Quelle fraction du réservoir plein cela représente-t-il ?
\end{exercice}\end{blocLivre}
\begin{blocLivre}{ex-6-15}\begin{exercice}[3]{Inventer un problème}
Inventer un problème dont la solution est le calcul « $\frac25$ de 30 », puis le résoudre.
\end{exercice}\end{blocLivre}

\clearpage
